\documentclass[12pt, a4paper]{article}

% ── Shared preamble (src/lecturenotes.sty) ────────────────────────────────────
%   Packages, colors, boxes, \sol/\soleq and math macros all live in
%   src/*.sty. \input@path points LaTeX two folders up to find them.
\makeatletter
\def\input@path{{../../}}
\makeatother
\usepackage{lecturenotes}

% ── Local macros ──────────────────────────────────────────────────────────────
\providecommand{\bH}{\mathbf{H}}

% ── Local TikZ styles ─────────────────────────────────────────────────────────
\usetikzlibrary{decorations.pathreplacing}
\usepackage{pgfplots}
\pgfplotsset{compat=1.18}
\tikzset{
  layerlabel/.style={font=\small\itshape, gray},
  opbox/.style={draw=orange!70!black, fill=notebg, thick, rounded corners=4pt,
      minimum height=0.8cm, minimum width=1.7cm, font=\small, align=center},
  convbox/.style={draw=mainblue, fill=boxbg, thick, rounded corners=3pt,
      minimum height=0.75cm, minimum width=1.6cm, font=\small, align=center},
  brace/.style={decorate, decoration={brace, amplitude=4pt}, thick, gray},
  frame/.style={draw=gray!70, fill=gray!8, thick},
  bbox/.style={line width=1.3pt},
}

% \cuboid[<label y>]{<x>}{<depth: channels>}{<height: spatial>}{<color>}{<label>}
%   pseudo-3D block for a (C, H, W) tensor, centred vertically on y = 0
\newcommand{\cuboid}[6][-1.75]{%
  \pgfmathsetmacro{\off}{#4*0.3}
  \filldraw[fill=#5!20, draw=#5!80!black] (#2,-#4/2) rectangle ++(#3,#4);
  \filldraw[fill=#5!35, draw=#5!80!black]
  (#2,#4/2) -- ++(\off,\off) -- ++(#3,0) -- ++(-\off,-\off) -- cycle;
  \filldraw[fill=#5!50, draw=#5!80!black]
  (#2+#3,-#4/2) -- ++(\off,\off) -- ++(0,#4) -- ++(-\off,-\off) -- cycle;
  \node[below, font=\scriptsize, align=center] at (#2+#3/2,#1) {#6};
}

% \chstack{<x>}{<y>}{<size>}{<colour list, back to front>}
%   a stack of square feature maps / filter slices, offset diagonally
\newcommand{\chstack}[4]{%
  \foreach \col [count=\i from 0] in {#4} {
      \filldraw[fill=\col!25, draw=\col!80!black]
      (#1+0.15*\i, #2-0.15*\i) rectangle ++(#3, #3);
    }%
}

% \bubble{<GFLOPs>}{<top-1>}{<radius pt>}{<colour>}{<label>}{<label anchor>}
%   one model in the bubble chart; radius = 1.6pt * sqrt(params in millions),
%   so the circle's area is proportional to the number of parameters
\newcommand{\bubble}[6]{%
  \fill[#4, fill opacity=0.35] (axis cs:#1,#2) circle[radius=#3pt];
  \draw[#4!80!black, thin] (axis cs:#1,#2) circle[radius=#3pt];
  \node[font=\scriptsize, anchor=#6, inner sep=2pt] at (axis cs:#1,#2) {#5};
}

% \scene{<nut1 style>}{<nut2 style>}{<bolt style>} — toy image: two nuts, one bolt
\newcommand{\scene}[3]{%
  \draw[frame] (0,0) rectangle (3.2,2.6);
  \filldraw[#1] (0.9,1.0) circle (0.45);
  \filldraw[#2] (2.3,1.85) circle (0.4);
  \filldraw[#3] (1.7,0.35) rectangle (2.9,0.75);
}

% ── Document ──────────────────────────────────────────────────────────────────
\begin{document}

\begin{center}
  {\LARGE \textbf{Section 4.2: CNN Architectures and Vision Tasks}} \\[0.3em]
  {\large\itshape From LeNet-5 to ConvNeXt, and Beyond Classification} \\[0.5em]
  {\large AI-IM Class --- Lecture Notes}%
  \ifstudentmode
    \quad {\normalsize\color{gray}[Student Copy]}%
  \fi
\end{center}

\vspace{1em}
\hrule
\vspace{1.5em}

With convolution, pooling and the classifier head (Section 4.1) we can build
any CNN. The question is \emph{which} CNN to build: how many layers, what
kernel sizes, how many channels. The answers came from a short series of
famous architectures, each introducing one idea that is still used today.
In this section we
\begin{itemize}[nosep]
  \item trace LeNet-5, the first successful CNN, layer by layer,
  \item follow the ImageNet competition and the key idea of each famous
        architecture, from AlexNet to ConvNeXt,
  \item compare them by accuracy, size and computing cost, and reuse a
        pretrained network for our own problem (\emph{transfer learning}),
        and
  \item look beyond classification: object detection, segmentation and OCR.
\end{itemize}

% ══════════════════════════════════════════════════════════════════════════════
\section{LeNet-5 (1998)}

\textbf{LeNet-5}, by Yann LeCun and colleagues, read handwritten digits on
bank cheques. It already has the structure of Section 4.1: two conv--pool
blocks followed by fully connected layers.

\begin{figure}[H]
  \centering
  \begin{tikzpicture}[>=Stealth, scale=0.95, every node/.style={transform shape}]
    \cuboid{0}{0.1}{3.2}{gray}{input\\$(1,32,32)$}
    \cuboid{1.9}{0.4}{2.8}{mainblue}{C1\\$(6,28,28)$}
    \cuboid{3.9}{0.4}{1.4}{mainblue}{S2\\$(6,14,14)$}
    \cuboid{5.6}{0.9}{1.0}{orange}{C3\\$(16,10,10)$}
    \cuboid{7.9}{0.9}{0.5}{orange}{S4\\$(16,5,5)$}
    \draw[->, thick] (1.1,0) -- (1.8,0);
    \draw[->, thick] (3.1,0) -- (3.8,0);
    \draw[->, thick] (4.8,0) -- (5.5,0);
    \draw[->, thick] (6.9,0) -- (7.8,0);

    \filldraw[fill=green!20, draw=green!50!black] (10.0,-1.2) rectangle ++(0.25,2.4);
    \node[below, font=\scriptsize, align=center] at (10.12,-1.75) {C5\\$120$};
    \filldraw[fill=green!20, draw=green!50!black] (11.3,-0.85) rectangle ++(0.25,1.7);
    \node[below, font=\scriptsize, align=center] at (11.42,-1.75) {F6\\$84$};
    \filldraw[fill=green!35, draw=green!50!black] (12.6,-0.4) rectangle ++(0.25,0.8);
    \node[below, font=\scriptsize, align=center] at (12.72,-1.75) {output\\$10$};
    \draw[->, thick] (9.1,0) -- (9.9,0);
    \draw[->, thick] (10.35,0) -- (11.2,0);
    \draw[->, thick] (11.65,0) -- (12.5,0);
  \end{tikzpicture}
  \caption{LeNet-5. C = convolution ($5\times5$, valid), S = subsampling
    ($2\times2$ average pooling), F = fully connected.}
  \label{fig:lenet}
\end{figure}

The input is a $32 \times 32$ grayscale image (MNIST digits padded from
$28 \times 28$). All convolutions use $5 \times 5$ kernels with stride $1$ and
no padding; all pooling layers are $2 \times 2$ with stride $2$. We use the
common modern version: ReLU activations, average pooling, every C3 filter
connected to all 6 input channels, and a linear output layer.
Table~\ref{tab:lenet} traces the shapes and parameters with the formulas of
Section 4.1.

\begin{table}[H]
  \centering
  \renewcommand{\arraystretch}{1.35}
  \begin{tabular}{llll}
    \toprule
    Layer                     & Operation              & Output shape    & Parameters                         \\
    \midrule
    Input                     &                        & $(1, 32, 32)$   & ---                                \\
    C1                        & 6 filters $5\times5$   & $(6, 28, 28)$   & $(1\cdot25+1)\cdot6 = 156$         \\
    S2                        & avg pool $2\times2$    & $(6, 14, 14)$   & $0$                                \\
    C3                        & 16 filters $5\times5$  & $(16, 10, 10)$  & $(6\cdot25+1)\cdot16 = 2\,416$     \\
    S4                        & avg pool $2\times2$    & $(16, 5, 5)$    & $0$                                \\
    C5                        & 120 filters $5\times5$ & $(120, 1, 1)$   & $(16\cdot25+1)\cdot120 = 48\,120$  \\
    F6                        & linear                 & $(84)$          & $(120+1)\cdot84 = 10\,164$         \\
    Output                    & linear                 & $(10)$          & $(84+1)\cdot10 = 850$              \\
    \midrule
    \multicolumn{3}{l}{Total} & $61\,706$                                                                   \\
    \bottomrule
  \end{tabular}
  \caption{LeNet-5 layer by layer.}
  \label{tab:lenet}
\end{table}

C5 is called a convolution although it behaves like a fully connected layer:
its $5\times5$ kernel exactly covers the $5\times5$ input, so there is only
one position and each filter sees every input value.

\begin{notebox}
  \textbf{The original LeNet-5} used scaled $\tanh$ activations, pooling
  layers with one trainable weight and bias per channel, and a hand-designed
  table that connected each C3 filter to only some of the S2 channels
  (1\,516 parameters instead of 2\,416), to save computation in 1998.
  Its total is still about 60 thousand parameters.
\end{notebox}

% ══════════════════════════════════════════════════════════════════════════════
\section{The ImageNet Challenge}

LeNet-5 worked on small grayscale digits. Natural photos are much harder:
thousands of object types, varying lighting, viewpoint and background.
The \textbf{ImageNet Large Scale Visual Recognition Challenge (ILSVRC)} ran
from 2010 to 2017. Its classification task had about 1.2 million training
images in 1\,000 classes. Performance is measured by the \textbf{top-5 error}:
a prediction counts as correct if the true class is among the model's five
most probable classes. The stricter \textbf{top-1 accuracy} counts only the
single most probable class. The jump in 2012 is the moment deep learning took
over computer vision; by 2015 CNNs beat the human estimate. Each of the
architectures below won (or nearly won) one of these years.

\begin{figure}[H]
  \centering
  \begin{tikzpicture}[x=1.45cm, y=0.13cm]
    \draw[->, gray] (0.3,0) -- (8.7,0);
    \draw[->, gray] (0.3,0) -- (0.3,32) node[above, font=\small] {top-5 error (\%)};
    \foreach \t in {0,10,20,30}
    \node[left, font=\scriptsize] at (0.3,\t) {\t};
    \draw[red!70!black, dashed, thick] (0.4,5.1) -- (8.6,5.1)
    node[right, font=\scriptsize, align=left] {human\\$\approx 5.1$};
    \foreach \i/\yr/\e/\name/\col in {
    1/2010/28.2/{}/gray,
    2/2011/25.8/{}/gray,
    3/2012/16.4/AlexNet/mainblue,
    4/2013/11.7/ZFNet/mainblue,
    5/2014/6.7/GoogLeNet/mainblue,
    6/2015/3.6/ResNet/mainblue,
    7/2016/3.0/{}/mainblue,
    8/2017/2.3/SENet/mainblue} {
    \filldraw[fill=\col!40, draw=\col!80!black] (\i-0.3,0) rectangle (\i+0.3,\e);
    \node[above, font=\scriptsize, fill=white, inner sep=1pt] at (\i,\e+0.4) {\e};
    \node[below, font=\scriptsize] at (\i,0) {\yr};
    \node[below, font=\scriptsize\itshape, gray] at (\i,-2.8) {\name};
    }
  \end{tikzpicture}
  \caption{Winning top-5 error in the ILSVRC classification task. Grey: methods
    based on hand-designed features; blue: CNNs. The 2016 winner combined
    existing networks rather than introducing a new one. The human estimate is
    Andrej Karpathy's result on a test sample.}
  \label{fig:ilsvrc}
\end{figure}

% ══════════════════════════════════════════════════════════════════════════════
\section{AlexNet (2012): Scale It Up}

\textbf{AlexNet} (Krizhevsky, Sutskever and Hinton) is LeNet made bigger:
5 convolutional and 3 fully connected layers, about 60 million parameters,
trained on $224 \times 224$ colour images. Its architecture was not new; what
made it work were practical choices that are still standard today:
\textbf{ReLU} instead of $\tanh$, which trained several times faster;
\textbf{GPUs} (two gaming graphics cards, for about a week);
\textbf{data augmentation} with random crops and flips of the training images;
and \textbf{dropout} in the fully connected layers against overfitting. The
2013 winner, \textbf{ZFNet} (Zeiler and Fergus), kept AlexNet's design but
used a smaller $7 \times 7$ first-layer kernel, a change they found by
\emph{visualising} what each layer had learned.

% ══════════════════════════════════════════════════════════════════════════════
\section{VGG (2014): Only Small Kernels}

\textbf{VGG} (Simonyan and Zisserman, Oxford) asked how far a very simple
design can go. It uses only $3 \times 3$ convolutions with stride $1$ and same
padding, and $2 \times 2$ max pooling, repeated in blocks until the network has
16 or 19 layers; each block halves the feature maps and doubles the channels
($64 \to 128 \to 256 \to 512$). Two stacked $3 \times 3$ layers have the same
receptive field as one $5 \times 5$ layer, but need $18C^2$ instead of $25C^2$
weights and add one more ReLU. VGG-16 reached 7.3\% top-5 error, but it has
about 138 million parameters, most of them in its three fully connected
layers. Its uniform design made it popular as a starting point for other
tasks.

% ══════════════════════════════════════════════════════════════════════════════
\section{GoogLeNet / Inception (2014): Go Wide, Cheaply}

\textbf{GoogLeNet} (Szegedy et al., Google) won 2014 with 6.7\% top-5 error
using about 7 million parameters, roughly 20 times fewer than VGG. Its
\textbf{Inception module} applies several kernel sizes in parallel and
concatenates the results along the channel dimension. What makes this
affordable is the \textbf{$1 \times 1$ convolution}: it looks at a single pixel
but across \emph{all} its channels, which is a small fully connected layer
applied at each pixel, so it changes the number of channels without touching
the height and width. Placing a $1\times1$ layer with few filters before each
$3\times3$ and $5\times5$ convolution (a \emph{bottleneck}) cuts their cost by
about an order of magnitude.

\begin{figure}[H]
  \centering
  \begin{tikzpicture}[>=Stealth, x=3.2cm, y=1.4cm]
    \node[opbox, minimum width=3cm] (in) at (1.5, 3.4) {previous layer};
    \node[convbox] (a1) at (0, 1.5) {$1\times1$};
    \node[convbox] (b1) at (1, 2) {$1\times1$};
    \node[convbox] (b2) at (1, 1) {$3\times3$};
    \node[convbox] (c1) at (2, 2) {$1\times1$};
    \node[convbox] (c2) at (2, 1) {$5\times5$};
    \node[opbox, minimum width=1.6cm] (d1) at (3, 2) {$3\times3$\\[-2pt]\scriptsize max pool};
    \node[convbox] (d2) at (3, 1) {$1\times1$};
    \node[opbox, minimum width=3cm] (out) at (1.5, -0.5) {concatenate channels};
    \draw[->] (in) -- (a1); \draw[->] (in) -- (b1);
    \draw[->] (in) -- (c1); \draw[->] (in) -- (d1);
    \draw[->] (b1) -- (b2); \draw[->] (c1) -- (c2); \draw[->] (d1) -- (d2);
    \draw[->] (a1) -- (out); \draw[->] (b2) -- (out);
    \draw[->] (c2) -- (out); \draw[->] (d2) -- (out);
  \end{tikzpicture}
  \caption{An Inception module. Four branches see the same input with
    different receptive fields. In the second and third branches a
    $1 \times 1$ convolution first shrinks the number of channels, so the
    expensive $3\times3$ and $5\times5$ convolutions become cheap.}
  \label{fig:inception}
\end{figure}

% ══════════════════════════════════════════════════════════════════════════════
\section{ResNet (2015): Skip Connections}

He et al.\ (Microsoft Research) found that a plain 56-layer network had a
\emph{higher training error} than a 20-layer one: deeper networks were harder
to \emph{optimise}. Their fix was the \textbf{residual block}, whose layers
compute a correction $F(\bH)$ that is \emph{added} to the block's input
through a \textbf{skip connection}: $\text{output} = \relu\bigl(F(\bH) +
\bH\bigr)$. If the best thing a block can do is nothing, it only needs
$F = 0$, which is easy to learn. The skip connection also keeps the gradient
alive: a plain layer $y = f(x)$ multiplies the gradient by $f'$, and a product
of many small factors vanishes (Section 3.4), while the residual version
$y = x + f(x)$ multiplies it by $1 + f'$, which does not shrink. With skip
connections ResNet trained networks with 152 layers and won 2015 with 3.6\%
top-5 error, better than the human estimate. Skip connections are now
everywhere in deep learning, including the transformers behind large language
models.

\begin{figure}[H]
  \centering
  \begin{tikzpicture}[>=Stealth, node distance=0.8cm]
    \node (in) {$\bH$};
    \node[convbox, right=of in] (c1) {$3\times3$ conv};
    \node[opbox, right=0.6cm of c1, minimum width=1.2cm] (r1) {ReLU};
    \node[convbox, right=0.6cm of r1] (c2) {$3\times3$ conv};
    \node[circle, draw, thick, right=1.3cm of c2, inner sep=1pt] (plus) {$+$};
    \node[opbox, right=0.6cm of plus, minimum width=1.2cm] (r2) {ReLU};
    \node[right=0.6cm of r2] (out) {output};
    \draw[->, thick] (in) -- (c1);
    \draw[->, thick] (c1) -- (r1);
    \draw[->, thick] (r1) -- (c2);
    \draw[->, thick] (c2) -- node[below, font=\small] {$F(\bH)$} (plus);
    \draw[->, thick] (plus) -- (r2);
    \draw[->, thick] (r2) -- (out);
    \draw[->, thick, red!70!black] ($(in.east)+(0.2,0)$) -- ++(0,-1.1)
    -| node[pos=0.25, below, font=\small] {identity (skip connection)} (plus);
    \draw[brace] ($(c1.north west)+(0,0.15)$) --
    node[above=5pt, font=\small] {$F$: the residual branch}
    ($(c2.north east)+(0,0.15)$);
  \end{tikzpicture}
  \caption{A residual block. The input skips past two convolutions and is
    added to their output. (Real ResNets also put a normalisation layer after
    each convolution, which we meet in a later topic.)}
  \label{fig:resblock}
\end{figure}

% ══════════════════════════════════════════════════════════════════════════════
\section{Efficient Networks (2016--2019): Small and Fast}

After ResNet, a new question became just as important as accuracy: can the
network run on a phone, a drone or a smart camera on a production line? Two
numbers measure the cost of a network: the number of \textbf{parameters},
which sets the memory needed to store the model, and the number of
\textbf{multiply--adds} per image, which sets the running time and energy and
is usually reported in \textbf{GFLOPs} (billions of operations).

% ──────────────────────────────────────────────────────────────────────────────
\subsection{SqueezeNet (2016): AlexNet Accuracy in 5 MB}

\textbf{SqueezeNet} (Iandola et al.) takes the $1\times1$ bottleneck of
GoogLeNet to the extreme. Its \emph{fire module} first \emph{squeezes} the
channels with a few $1\times1$ filters, then \emph{expands} them again with a
mix of $1\times1$ and $3\times3$ filters, and the network has no large fully
connected layers. It matches AlexNet's accuracy with about 1.2 million
parameters, roughly 50 times fewer, so the whole model fits in less than
5\,MB.

% ──────────────────────────────────────────────────────────────────────────────
\subsection{MobileNet (2017): Depthwise Separable Convolutions}

A standard convolution filters space and mixes all channels at once.
\textbf{MobileNet} (Howard et al., Google) splits these two jobs into a
\textbf{depthwise} convolution (one $k \times k$ filter per input channel,
applied to that channel only) followed by a \textbf{pointwise} $1 \times 1$
convolution that mixes the channels. This \textbf{depthwise separable
convolution} costs $1/C_{\text{out}} + 1/k^2$ of a standard one: for
$3 \times 3$ kernels and many channels, about 8 to 9 times fewer weights and
multiplications, for a small loss in accuracy. \textbf{MobileNetV2} (2018)
adds ResNet's skip connections, and \textbf{MobileNetV3} (2019) reaches 74.0\%
top-1 accuracy with almost 20 times less computation than ResNet-50.

\begin{figure}[H]
  \centering
  \begin{tikzpicture}[>=Stealth, font=\small]
    % ── (a) standard convolution ─────────────────────────────────────────────
    % labels hang from y = -0.8, below the lowest drawing (y = -0.5)
    \tikzset{steplabel/.style={layerlabel, anchor=north, align=center}}
    \node[anchor=west, font=\small\bfseries] at (-0.2, 2.2) {(a) Standard convolution};
    \chstack{0}{0}{1.4}{orange, mainblue, green!60!black}
    \node[steplabel] at (0.85,-0.8) {input\\$C_{\text{in}} = 3$};
    \draw[->, thick] (2.0,0.35) -- (3.3,0.35);
    \chstack{3.6}{0.2}{0.6}{orange, mainblue, green!60!black}
    \node[font=\small] at (5.0,0.35) {$\times\,4$};
    \node[steplabel] at (4.25,-0.8) {4 filters, each\\$k \times k \times C_{\text{in}}$};
    \draw[->, thick] (5.5,0.35) -- (6.8,0.35);
    \chstack{7.1}{0}{1.4}{gray, gray, gray, gray}
    \node[steplabel] at (8.0,-0.8) {output\\$C_{\text{out}} = 4$};
    \node[anchor=west, align=left] at (10.0,0.35) {weights:\\$k^2\, C_{\text{in}}\, C_{\text{out}}$};

    % ── (b) depthwise separable convolution ──────────────────────────────────
    \begin{scope}[yshift=-4.8cm]
      \node[anchor=west, font=\small\bfseries] at (-0.2, 2.2) {(b) Depthwise separable convolution};
      \chstack{0}{0}{1.4}{orange, mainblue, green!60!black}
      \node[steplabel] at (0.85,-0.8) {input\\$C_{\text{in}} = 3$};
      \draw[->, thick] (2.0,0.35) -- (3.0,0.35);
      \filldraw[fill=orange!25, draw=orange!80!black] (3.3,1.0) rectangle ++(0.5,0.5);
      \filldraw[fill=mainblue!25, draw=mainblue!80!black] (3.3,0.25) rectangle ++(0.5,0.5);
      \filldraw[fill=green!60!black!25, draw=green!60!black!80!black] (3.3,-0.5) rectangle ++(0.5,0.5);
      \node[steplabel] at (3.55,-0.8) {depthwise:\\one $k\times k$\\per channel};
      \draw[->, thick] (4.1,0.35) -- (5.1,0.35);
      \chstack{5.4}{0}{1.4}{orange, mainblue, green!60!black}
      \node[steplabel] at (6.25,-0.8) {still\\$C_{\text{in}}$ maps};
      \draw[->, thick] (7.4,0.35) -- (8.4,0.35);
      \chstack{8.7}{0.3}{0.3}{orange, mainblue, green!60!black}
      \node[font=\small] at (9.75,0.35) {$\times\,4$};
      \node[steplabel] at (9.1,-0.8) {pointwise:\\4 filters, each\\$1\times1\times C_{\text{in}}$};
      \draw[->, thick] (10.2,0.35) -- (11.2,0.35);
      \chstack{11.5}{0}{1.4}{gray, gray, gray, gray}
      \node[steplabel] at (12.4,-0.8) {output\\$C_{\text{out}} = 4$};
      \node[anchor=north] at (6.7,-2.5)
      {weights: $\underbrace{k^2\, C_{\text{in}}}_{\text{depthwise}}
          + \underbrace{C_{\text{in}}\, C_{\text{out}}}_{\text{pointwise}}$};
    \end{scope}
  \end{tikzpicture}
  \caption{(a) Each standard filter covers $k \times k$ pixels in all input
    channels at once. (b) MobileNet first filters each channel separately with
    its own $k \times k$ filter (the colours stay separate), then mixes the
    channels with $1 \times 1$ filters.}
  \label{fig:depthwise}
\end{figure}

% ──────────────────────────────────────────────────────────────────────────────
\subsection{EfficientNet (2019): Scale Everything Together}

Given more computing power, should a network grow deeper (more layers), wider
(more channels), or take larger input images? \textbf{EfficientNet} (Tan and
Le, Google) grows all three together by fixed factors, which they call
\textbf{compound scaling}. The smallest model, EfficientNet-B0, is built from
MobileNet-style blocks; B1 to B7 are scaled-up copies of it. EfficientNet-B0
beats ResNet-50 with one fifth of the parameters and one tenth of the
computation, and B7 reaches 84.1\%.

% ──────────────────────────────────────────────────────────────────────────────
\subsection{Squeeze-and-Excitation (SENet, 2017)}

The last ILSVRC winner (Hu et al., 2.3\% top-5 error) is not a new
architecture but an add-on block for ResNet or Inception that lets the network
weight its channels for each image. It \emph{squeezes} every channel to one
number with global average pooling, passes this vector through a tiny MLP
ending in a sigmoid, and multiplies each channel by the resulting weight
between 0 and 1, so a ``fur'' channel can be turned up for a dog and down for
a car. MobileNetV3 and EfficientNet use such blocks throughout, and this
input-dependent reweighting is an early form of the \emph{attention} used by
transformers.

% ══════════════════════════════════════════════════════════════════════════════
\section{After 2020: Transformers and Modern CNNs}

\paragraph{Vision Transformer (ViT, 2020).}
Dosovitskiy et al.\ (Google) removed convolution altogether. A ViT cuts the
image into patches of $16 \times 16$ pixels, flattens each patch into a
vector and treats the patches like the words of a sentence. A
\textbf{transformer}, the same architecture as in large language models,
processes this sequence; its \textbf{attention} layers let every patch look at
every other patch directly. The price is that a ViT has no built-in idea that
nearby pixels belong together: it must learn this from data. Trained on
ImageNet alone it is no better than a good CNN; pretrained on hundreds of
millions of images it gives the best results.

\paragraph{ConvNeXt (2022).}
Liu et al.\ (Meta) started from ResNet-50 and modernised it step by step:
larger $7\times7$ depthwise kernels, fewer activation functions, and the long
training with heavy augmentation used for transformers. The resulting pure
CNN, \textbf{ConvNeXt}, matches or beats ViT of the same size. Much of the
progress came from training recipes rather than from attention itself:
simply retraining the \emph{original} ResNet-50 with a modern recipe raises
its top-1 accuracy from 76.1\% to 80.9\%.

% ══════════════════════════════════════════════════════════════════════════════
\section{Comparison}
\label{sec:comparison}

Table~\ref{tab:comparison} and Figure~\ref{fig:bubble} compare the networks of
this section. All numbers are for the pretrained models in
\texttt{torchvision}, the library we use for transfer learning below:
top-1 accuracy on the ImageNet validation set with a single
$224 \times 224$ input, and GFLOPs counted as multiply--adds per image. They
differ slightly from the competition results quoted earlier.

\begin{table}[H]
  \centering
  \small
  \renewcommand{\arraystretch}{1.3}
  \begin{tabular}{lcrrrl}
    \toprule
    Network           & Year & Params   & GFLOPs & Top-1  & Key idea                         \\
    \midrule
    LeNet-5           & 1998 & 0.06\,M  & ---    & ---    & conv + pool + dense              \\
    AlexNet           & 2012 & 61.1\,M  & 0.71   & 56.5\% & ReLU, GPUs, augmentation         \\
    VGG-16            & 2014 & 138.4\,M & 15.5   & 71.6\% & stacks of $3\times3$ kernels     \\
    GoogLeNet         & 2014 & 6.6\,M   & 1.50   & 69.8\% & Inception, $1\times1$ bottleneck \\
    ResNet-18         & 2015 & 11.7\,M  & 1.81   & 69.8\% & skip connections                 \\
    ResNet-50         & 2015 & 25.6\,M  & 4.09   & 76.1\% & skip connections                 \\
    ResNet-152        & 2015 & 60.2\,M  & 11.5   & 78.3\% & skip connections                 \\
    SqueezeNet 1.1    & 2016 & 1.2\,M   & 0.35   & 58.2\% & fire modules                     \\
    MobileNetV2       & 2018 & 3.5\,M   & 0.30   & 71.9\% & separable conv + skips           \\
    MobileNetV3-Large & 2019 & 5.5\,M   & 0.22   & 74.0\% & automated search                 \\
    EfficientNet-B0   & 2019 & 5.3\,M   & 0.39   & 77.7\% & compound scaling                 \\
    EfficientNet-B7   & 2019 & 66.3\,M  & 37.7   & 84.1\% & compound scaling                 \\
    ViT-B/16          & 2020 & 86.6\,M  & 17.6   & 81.1\% & transformer on patches           \\
    ConvNeXt-Tiny     & 2022 & 28.6\,M  & 4.46   & 82.5\% & modernised ResNet                \\
    \bottomrule
  \end{tabular}
  \caption{Famous image classifiers. Accuracy, parameters and GFLOPs are
    those of the \texttt{IMAGENET1K\_V1} weights in \texttt{torchvision}.}
  \label{tab:comparison}
\end{table}

\begin{figure}[H]
  \centering
  \begin{tikzpicture}
    \begin{axis}[
        width=15cm, height=9.5cm,
        xmode=log, xmin=0.1, xmax=80, ymin=50, ymax=88,
        log ticks with fixed point,
        xtick={0.1,0.3,1,3,10,30},
        xlabel={computation per image (GFLOPs, log scale)},
        ylabel={ImageNet top-1 accuracy (\%)},
        grid=major, grid style={gray!20},
        tick label style={font=\small}, label style={font=\small},
        clip=false,
      ]
      % classic CNNs (2012–2015)
      \bubble{0.714}{56.5}{12.5}{mainblue}{\hspace*{13pt}AlexNet}{west}
      \bubble{15.47}{71.6}{18.8}{mainblue}{VGG-16}{center}
      \bubble{1.498}{69.8}{4.1}{mainblue}{GoogLeNet\hspace*{7pt}}{east}
      \bubble{1.814}{69.8}{5.5}{mainblue}{\hspace*{8pt}ResNet-18}{west}
      \bubble{4.089}{76.1}{8.1}{mainblue}{\hspace*{10pt}ResNet-50}{west}
      \bubble{11.51}{78.3}{12.4}{mainblue}{ResNet-152\hspace*{14pt}}{east}
      % efficient CNNs (2016–2019)
      \bubble{0.349}{58.2}{2.0}{green!60!black}{SqueezeNet\hspace*{3pt}}{east}
      \bubble{0.301}{71.9}{3.0}{green!60!black}{\hspace*{5pt}MobileNetV2}{west}
      \bubble{0.217}{74.0}{3.7}{green!60!black}{\hspace*{5pt}MobileNetV3-L}{west}
      \bubble{0.386}{77.7}{3.7}{green!60!black}{\hspace*{5pt}EfficientNet-B0}{west}
      \bubble{37.75}{84.1}{13.0}{green!60!black}{EfficientNet-B7\hspace*{15pt}}{east}
      % 2020 and later
      \bubble{17.56}{81.1}{14.9}{orange}{ViT-B/16\hspace*{16pt}}{east}
      \bubble{4.456}{82.5}{8.6}{orange}{ConvNeXt-T\hspace*{10pt}}{east}

      % legend (lower right)
      \node[anchor=west, font=\scriptsize] at (axis cs:4,66) {%
        \tikz\fill[mainblue, fill opacity=0.5] (0,0) circle (3pt); classic CNN (2012--15)};
      \node[anchor=west, font=\scriptsize] at (axis cs:4,63.8) {%
        \tikz\fill[green!60!black, fill opacity=0.5] (0,0) circle (3pt); efficient CNN (2016--19)};
      \node[anchor=west, font=\scriptsize] at (axis cs:4,61.6) {%
        \tikz\fill[orange, fill opacity=0.5] (0,0) circle (3pt); 2020 and later};
      \draw[gray] (axis cs:9,56.6)  circle[radius=3.6pt];
      \draw[gray] (axis cs:17,56.6) circle[radius=8pt];
      \draw[gray] (axis cs:36,56.6) circle[radius=16pt];
      \node[font=\scriptsize, gray] at (axis cs:9,52.2)  {5\,M};
      \node[font=\scriptsize, gray] at (axis cs:17,52.2) {25\,M};
      \node[font=\scriptsize, gray] at (axis cs:36,52.2) {100\,M};
      \node[font=\scriptsize, gray, anchor=east] at (axis cs:7,56.6) {parameters:};
    \end{axis}
  \end{tikzpicture}
  \caption{Accuracy against computation; the area of each circle is
    proportional to the number of parameters. Good models sit in the upper
    left: high accuracy for little computation. Data from
    Table~\ref{tab:comparison}.}
  \label{fig:bubble}
\end{figure}

% ══════════════════════════════════════════════════════════════════════════════
\section{Transfer Learning}
\label{sec:transfer}

Training a ResNet from scratch needs a million labelled images and days of
GPU time, but a factory that wants to classify, say, three kinds of surface
defect may have only a few hundred photos. \textbf{Transfer learning} reuses a
network that was already trained on ImageNet. Its feature extractor has
learned general features (edges, textures, shapes) that are useful for almost
any image; only the classifier head is specific to the 1\,000 ImageNet
classes. So we load the pretrained network, replace its final linear layer
with a new one that has one output per \emph{our} class, and train on our
data: either only the new layer (\emph{feature extraction}, with the rest
frozen), or the whole network with a small learning rate
(\emph{fine-tuning}).

\begin{codebox}
import torchvision

model = torchvision.models.resnet18(weights="IMAGENET1K_V1")
for p in model.parameters():
    p.requires_grad = False              # freeze the feature extractor

model.fc = nn.Linear(512, 3)             # new head: 3 defect classes
optimizer = torch.optim.Adam(model.fc.parameters(), lr=1e-3)
\end{codebox}

% ══════════════════════════════════════════════════════════════════════════════
\section{Beyond Classification: A Family of Vision Tasks}

\begin{figure}[H]
  \centering
  \begin{tikzpicture}[scale=0.95]
    % (a) classification
    \begin{scope}
      \draw[frame] (0,0) rectangle (3.2,2.6);
      \filldraw[fill=gray!45, draw=gray!80!black] (1.6,1.3) circle (0.7);
      \filldraw[fill=gray!8, draw=gray!80!black] (1.6,1.3) circle (0.3);
      \node[font=\small] at (1.6,-0.35) {label: \texttt{nut}};
      \node[font=\small\bfseries] at (1.6,-0.95) {(a) classification};
    \end{scope}
    % (b) detection
    \begin{scope}[xshift=3.7cm]
      \scene{fill=gray!45, draw=gray!80!black}{fill=gray!45, draw=gray!80!black}
      {fill=gray!45, draw=gray!80!black}
      \draw[bbox, mainblue] (0.4,0.5) rectangle (1.4,1.5);
      \draw[bbox, mainblue] (1.85,1.4) rectangle (2.75,2.3);
      \draw[bbox, orange!80!black] (1.6,0.25) rectangle (3.0,0.85);
      \node[font=\tiny, fill=mainblue, text=white, inner sep=1pt, anchor=south west]
      at (0.4,1.5) {nut 0.97};
      \node[font=\tiny, fill=mainblue, text=white, inner sep=1pt, anchor=north east]
      at (1.85,2.3) {nut 0.91};
      \node[font=\tiny, fill=orange!80!black, text=white, inner sep=1pt, anchor=north west]
      at (1.6,0.25) {bolt 0.88};
      \node[font=\small] at (1.6,-0.35) {boxes + labels};
      \node[font=\small\bfseries] at (1.6,-0.95) {(b) object detection};
    \end{scope}
    % (c) semantic segmentation
    \begin{scope}[xshift=7.4cm]
      \scene{fill=mainblue!60, draw=mainblue}{fill=mainblue!60, draw=mainblue}
      {fill=orange!60, draw=orange!80!black}
      \node[font=\small] at (1.6,-0.35) {class of every pixel};
      \node[font=\small\bfseries] at (1.6,-0.95) {(c) semantic segm.};
    \end{scope}
    % (d) instance segmentation
    \begin{scope}[xshift=11.1cm]
      \scene{fill=mainblue!60, draw=mainblue}{fill=green!50!black!50, draw=green!40!black}
      {fill=orange!60, draw=orange!80!black}
      \node[font=\small] at (1.6,-0.35) {pixels of every object};
      \node[font=\small\bfseries] at (1.6,-0.95) {(d) instance segm.};
    \end{scope}
  \end{tikzpicture}
  \caption{Four vision tasks on images of parts. (a) One label for the whole
    image. (b) A box, class and confidence score for every object. (c) Each
    pixel is coloured by class; the two nuts get the same colour. (d) Each
    object gets its own mask, so the two nuts are distinguished.}
  \label{fig:tasks}
\end{figure}

\begin{center}
  \renewcommand{\arraystretch}{1.35}
  \begin{tabular}{lll}
    \toprule
    Task                   & Output per image                        & Loss (main parts)     \\
    \midrule
    Classification         & $C$ class scores                        & cross-entropy         \\
    Object detection       & a list of (box, class, score)           & box regression + CE   \\
    Semantic segmentation  & $C$ scores per pixel, shape $(C, H, W)$ & CE at every pixel     \\
    Instance segmentation  & boxes + a pixel mask per object         & detection + mask loss \\
    \bottomrule
  \end{tabular}
\end{center}

So far a CNN has answered one question per image: \emph{what is in it?} Many
real problems need more: a robot must know \emph{where} each part lies, an
inspection system must outline the scratch, and a traceability system must
\emph{read} the serial number. The same convolutional feature extractor
handles all of these; only the head and the loss change. Every loss is built
from the two pieces of Section 3.4: a \emph{regression} loss for continuous
outputs such as box coordinates, and \emph{cross-entropy} for class labels.
This part is an overview; the details are left to an advanced topic.

% ══════════════════════════════════════════════════════════════════════════════
\section{Object Detection}

A detector outputs, for every object, a \textbf{bounding box} (the smallest
axis-aligned rectangle around it, given by its corners or by its centre and
size), a class and a confidence score. A predicted box will never match the
true box exactly, so overlap is measured by the \textbf{intersection over
union}, $\text{IoU}(A, B) = \text{area}(A \cap B) / \text{area}(A \cup B)$,
which is $1$ for identical boxes and $0$ for boxes that do not touch; a
detection usually counts as correct if $\text{IoU} \ge 0.5$ with a true box of
the same class. Because a detector looks at many positions, it often reports
the same object several times. \textbf{Non-maximum suppression (NMS)} keeps
the highest-scoring box and deletes every other box of the same class whose
IoU with it is above a threshold, then repeats with the boxes that are left.

\begin{figure}[H]
  \centering
  \begin{tikzpicture}[x=0.7cm, y=0.7cm]
    \draw[gray!30, step=1] (0,0) grid (7,6);
    \foreach \x in {0,...,7} \node[below, font=\scriptsize, gray] at (\x,0) {\x};
    \foreach \y in {0,...,6} \node[left, font=\scriptsize, gray] at (0,\y) {\y};
    \fill[purple!30] (2,1) rectangle (4,4);
    \draw[bbox, mainblue] (0,0) rectangle (4,4);
    \draw[bbox, orange!80!black] (2,1) rectangle (6,5);
    \node[font=\small, mainblue] at (1,3.5) {$A$};
    \node[font=\small, orange!80!black] at (5.4,4.5) {$B$};
    \node[font=\scriptsize, purple!70!black] at (3,2.5) {$A\cap B$};
  \end{tikzpicture}
  \caption{Boxes $A = (0, 0, 4, 4)$ and $B = (2, 1, 6, 5)$ in corner format.
    Their intersection is shaded.}
  \label{fig:iou}
\end{figure}

\textbf{Two-stage detectors} (R-CNN, Faster R-CNN) first propose regions
that may contain an object, then classify each one; they are accurate but
slower. \textbf{One-stage detectors} such as \textbf{YOLO} (``You Only Look
Once'') predict all boxes in a single forward pass, fast enough for live
video, and are the usual first choice in industry. YOLO divides the image into
an $S \times S$ grid; the cell containing an object's centre predicts $B$
boxes (each with $x_c, y_c, w, h$ and a confidence) and $C$ class
probabilities, so the whole output is one tensor of shape
$S \times S \times (5B + C)$, followed by NMS. Detectors are evaluated by
\textbf{precision} (how many reported boxes are right) and \textbf{recall}
(how many true objects were found), summarised over all classes as the
\textbf{mAP}.

\begin{figure}[H]
  \centering
  \begin{tikzpicture}[x=0.75cm, y=0.75cm]
    \draw[frame] (0,0) rectangle (7,7);
    \fill[mainblue!15] (2,3) rectangle (3,4);
    \draw[gray!60, step=1] (0,0) grid (7,7);
    \filldraw[fill=gray!45, draw=gray!80!black] (2.6,3.4) circle (1.1);
    \filldraw[fill=gray!8, draw=gray!80!black] (2.6,3.4) circle (0.45);
    \fill[mainblue!45, opacity=0.55] (2,3) rectangle (3,4);
    \draw[mainblue, thick] (2,3) rectangle (3,4);
    \draw[bbox, mainblue] (1.5,2.3) rectangle (3.7,4.5);
    \fill[red!70!black] (2.6,3.4) circle (2.5pt);
    \node[font=\scriptsize, anchor=west, align=left] at (7.4,3.4)
    {the cell containing the\\object's centre (red dot)\\predicts its box};
    \draw[brace] (0,7.15) -- node[above=4pt, font=\small] {$S = 7$} (7,7.15);
  \end{tikzpicture}
  \caption{YOLO divides the image into an $S \times S$ grid. The shaded cell
    contains the centre of the nut and is responsible for its box. The
    object's centre can lie anywhere inside the cell; the cell predicts the
    box centre as an offset within it, and the box itself can be much larger
    than the cell.}
  \label{fig:yolo}
\end{figure}

% ══════════════════════════════════════════════════════════════════════════════
\section{Segmentation}

\textbf{Semantic segmentation} classifies every pixel: for $C$ classes the
network outputs a $(C, H, W)$ tensor of logits, with a softmax at each pixel
and the cross-entropy averaged over all pixels. A classification CNN shrinks
the feature maps, but segmentation needs the answer back at full resolution,
so it uses an \textbf{encoder--decoder}: an ordinary CNN encoder downsamples,
and a decoder upsamples again. The popular \textbf{U-Net} (2015) adds skip
connections that copy fine detail from each encoder level to the decoder level
of the same size. \textbf{Instance segmentation} (e.g.\ Mask R-CNN) combines
both ideas: a detector finds each object, and a small segmentation head
predicts a mask inside each box.

\begin{figure}[H]
  \centering
  \begin{tikzpicture}[>=Stealth, scale=0.95, every node/.style={transform shape}]
    \cuboid[-1.6]{0}{0.3}{2.8}{gray}{input\\$(3,H,W)$}
    \cuboid[-1.6]{1.6}{0.5}{2.2}{mainblue}{}
    \cuboid[-1.6]{3.3}{0.8}{1.4}{mainblue}{}
    \cuboid[-1.6]{5.4}{1.1}{0.8}{mainblue}{bottleneck}
    \cuboid[-1.6]{7.8}{0.8}{1.4}{orange}{}
    \cuboid[-1.6]{9.8}{0.5}{2.2}{orange}{}
    \cuboid[-1.6]{11.5}{0.3}{2.8}{green!60!black}{output\\$(C,H,W)$}
    \draw[->, thick] (1.0,0) -- (1.5,0);
    \draw[->, thick] (2.5,0) -- (3.2,0);
    \draw[->, thick] (4.4,0) -- (5.3,0);
    \draw[->, thick] (6.8,0) -- (7.7,0);
    \draw[->, thick] (9.0,0) -- (9.7,0);
    \draw[->, thick] (10.8,0) -- (11.4,0);
    \draw[->, red!70!black, thick, dashed] (2.0,1.7) to[bend left=18] (10.0,1.7);
    \draw[->, red!70!black, thick, dashed] (3.7,1.2) to[bend left=18] (8.2,1.2);
    \node[font=\small, red!70!black] at (6.0,2.9) {skip connections};
    \draw[brace] (4.9,-2.6) -- node[below=5pt, font=\small] {encoder (downsample)} (0,-2.6);
    \draw[brace] (12.4,-2.6) -- node[below=5pt, font=\small] {decoder (upsample)} (7.6,-2.6);
  \end{tikzpicture}
  \caption{An encoder--decoder for segmentation in the style of U-Net. The
    output has one score map per class, at the input resolution.}
  \label{fig:unet}
\end{figure}

% ══════════════════════════════════════════════════════════════════════════════
\section{Optical Character Recognition (OCR)}

\textbf{OCR} turns text in an image into a string: a serial number on a part,
a lot code and expiry date on a package, the reading of a meter. Modern OCR
systems are a pipeline of the tasks above. \textbf{Text detection} finds the
regions that contain text, using detection or segmentation. \textbf{Text
recognition} then reads each cropped line: instead of cutting it into single
characters, modern recognisers such as \textbf{CRNN} run a CNN over the whole
line, read the feature map column by column as a sequence, and are trained
with the \textbf{CTC} loss, which does not need to know where each character
is. Finally, \textbf{post-processing} corrects the string with known formats,
dictionaries or checksums. Ready-made libraries include \textbf{Tesseract},
\textbf{EasyOCR} and \textbf{PaddleOCR}.

\begin{figure}[H]
  \centering
  \begin{tikzpicture}[>=Stealth, node distance=0.75cm,
      opbox/.append style={minimum height=0.9cm, minimum width=1.9cm}]
    \node[opbox] (img) {image};
    \node[opbox, right=of img] (det) {text\\detection};
    \node[opbox, right=of det] (crop) {crop each\\text line};
    \node[opbox, right=of crop] (rec) {text\\recognition};
    \node[opbox, right=of rec] (post) {post-\\processing};
    \draw[->, thick] (img) -- (det);
    \draw[->, thick] (det) -- (crop);
    \draw[->, thick] (crop) -- (rec);
    \draw[->, thick] (rec) -- (post);
    \node[below=0.35cm of rec, font=\scriptsize\ttfamily] {L0T 2407-B};
    \node[below=0.35cm of post, font=\scriptsize\ttfamily] {LOT 2407-B};
    \node[below=0.35cm of det, font=\scriptsize\itshape, gray, align=center]
    {a detection or\\segmentation CNN};
  \end{tikzpicture}
  \caption{An OCR pipeline. Post-processing applies known rules (here: the
    code starts with the word \texttt{LOT}, so the digit \texttt{0} must be the
    letter \texttt{O}).}
  \label{fig:ocr}
\end{figure}

% ══════════════════════════════════════════════════════════════════════════════
\section{Vision in Manufacturing}

\begin{center}
  \renewcommand{\arraystretch}{1.35}
  \begin{tabular}{>{\raggedright\arraybackslash}p{5.2cm}
      >{\raggedright\arraybackslash}p{3.6cm}
      >{\raggedright\arraybackslash}p{5.4cm}}
    \toprule
    Application                         & Task                    & Example                                         \\
    \midrule
    Accept / reject a part              & classification          & surface defect type on steel sheet              \\
    Locate defects or count parts       & object detection        & missing components on a PCB; parts on a conveyor \\
    Measure defect shape and size       & segmentation            & crack or weld-seam outline                      \\
    Robot picking                       & detection / instance segm. & grasp points in a bin of parts              \\
    Traceability                        & OCR                     & serial numbers, lot codes, expiry dates         \\
    Rare or unknown defects             & anomaly detection       & train only on good parts, flag anything unusual \\
    \bottomrule
  \end{tabular}
\end{center}

\begin{notebox}
  \textbf{Practical advice.} Start from a pretrained model and fine-tune it
  (transfer learning, Section~\ref{sec:transfer}); a few hundred well-labelled
  images per class are often enough. Invest in consistent lighting and camera
  position: a controlled image makes every task easier. Choose the evaluation
  measure with the cost of errors in mind: a missed defect (low recall) is
  usually far more expensive than a false alarm.
\end{notebox}

% ══════════════════════════════════════════════════════════════════════════════
\section{Summary}

\begin{itemize}
  \item LeNet-5 set the pattern: conv--pool blocks followed by dense layers.
  \item AlexNet (2012) showed that a large CNN trained on GPUs with ReLU,
        augmentation and dropout beats hand-designed features by a wide
        margin.
  \item VGG: stacks of $3\times3$ kernels. GoogLeNet: parallel branches with
        $1\times1$ bottlenecks. ResNet: skip connections,
        $\text{output} = \relu(F(\bH) + \bH)$, make 100+ layers trainable.
  \item SqueezeNet, MobileNet (depthwise separable convolutions) and
        EfficientNet (compound scaling) reach good accuracy with few
        parameters and little computation.
  \item Vision transformers replace convolution with attention on image
        patches; ConvNeXt shows a modernised CNN does as well. Compare
        networks by top-1 accuracy, parameters and GFLOPs together.
  \item Transfer learning reuses an ImageNet feature extractor and retrains
        only a new head, so a few hundred images can be enough.
  \item Detection, segmentation and OCR share the CNN feature extractor and
        differ in the head and the loss. Detection uses boxes, IoU and NMS;
        YOLO predicts all boxes in one pass. Segmentation predicts a class per
        pixel with an encoder--decoder such as U-Net. OCR detects text, then
        recognises each line as a sequence.
\end{itemize}

\paragraph{Advanced topics.}
The YOLO loss function and anchor boxes, the details of CTC, transposed
convolutions and modern segmentation models, and anomaly detection will be
covered in a later topic.

\end{document}
